%! Expressions and Operations — Unit 1, Lesson 1.1 — Student Packet Cover Sheet
\documentclass[10pt]{article}
\usepackage{algebra2-article}
\usepackage{algebra2-boxes}
\usepackage{ltablex}
\keepXColumns

\begin{document}

% Full-bleed forest banner
\begin{tikzpicture}[remember picture, overlay]
  \fill[forest] (current page.north west)
    rectangle ([yshift=-0.9in]current page.north east);
\end{tikzpicture}
\vspace{-0.8in}

\begin{center}
  {\color{white}\textbf{\LARGE Algebra 2: Shepherd}}\\[4pt]
  {\color{forestmid}\large\textbf{Unit 1 \quad Foundations: Algebra 1 Review}}\\[6pt]
  {\color{white}\textbf{\large Lesson 1.1 \quad Expressions and Operations}}
\end{center}

\vspace{0.22in}
\namedateperiod
\vspace{0.18in}

\begin{learningtargetbox}
\small
\begin{enumerate}[leftmargin=1.5em, itemsep=5pt]
  \item \ldots{} \textbf{translate} a verbal quantitative situation into an algebraic expression --- and read one back into words --- without falling for the ``less than'' or ``the quantity'' traps.
  \item \ldots{} \textbf{evaluate} an expression at given replacement values, including absolute values, square roots, and cube roots, and explain why $-x^2$ and $(-x)^2$ are not the same thing.
  \item \ldots{} \textbf{derive} the laws of exponents by counting factors, then apply them to multivariable expressions and monomial ratios, leaving no negative exponent behind.
  \item \ldots{} add, subtract, multiply, divide, and \textbf{factor completely} polynomials in one variable --- GCF first, always --- recognizing the difference-of-squares and perfect-square patterns on sight and checking every factorization by expanding it.
  \item \ldots{} \textbf{simplify radicals} and move freely between $\sqrt[n]{a}$ and $a^{1/n}$.
  \item \ldots{} \textbf{justify} that two expressions are equivalent with algebra, and disprove it with a single replacement value.
\end{enumerate}
\end{learningtargetbox}

\vspace{0.14in}

\begin{tocbox}
\small
\renewcommand{\arraystretch}{1.5}
\begin{tabularx}{\linewidth}{@{} c l X r @{}}
\rowcolor{forestbg}
\textbf{\#} & \textbf{Component} & \textbf{Description} & \textbf{Score} \\
\midrule
1 & Warm-Up        & Order of operations, the $-x^2$ argument, and exponents re-derived by counting factors & \blank{1.2cm} \\
2 & Guided Notes   & Translating, evaluating, the exponent laws, polynomial operations and complete factoring, radicals & \blank{1.2cm} \\
3 & Group Activity & \emph{Rewriting Without Changing the Value} --- practice, mixed sets, error hunting, and structure, by tier & \blank{1.2cm} \\
4 & Exit Ticket    & Quick check: one evaluation, one complete factoring with a check, and an SOL-style exponent item & \blank{1.2cm} \\
5 & Homework       & The full domain area --- plus the equivalence audit & \blank{1.2cm} \\
\midrule
  & \multicolumn{2}{r}{\textbf{Total}} & \blank{1.2cm} \\
\end{tabularx}
\end{tocbox}

\vspace{0.10in}

\begin{remindbox}
\small An \textbf{expression} is a recipe for a number: it has no equals sign, so it is never
\emph{solved} --- it is \textbf{evaluated}, and only once you supply values. Everything in this
lesson is a rule for rewriting one expression as an \textbf{equivalent} one, meaning the two give
the same value for \emph{every} legal replacement value.
\\[3pt]
The \textbf{laws of exponents} are not facts to memorize. Write the factors out and count:
multiplying puts two piles together, so the counts add; dividing cancels in pairs, so they subtract;
a power of a power repeats a pile, so they multiply. Push the quotient law past where it ``should''
stop and it \emph{forces} $a^0=1$ and $a^{-n}=\frac{1}{a^n}$ --- neither is a convention somebody
invented. The same idea reaches radicals through $a^{1/2}=\sqrt{a}$ and $a^{1/3}=\sqrt[3]{a}$.
\\[3pt]
\textbf{Factoring is multiplying backwards}, which means every factorization can be checked by
expanding it --- so you never have to ask whether yours is right. Do it in order: \textbf{GCF
first, always}, then the pattern or the trinomial method, then check whether any factor \emph{still}
factors. That last check is the entire difference between ``factored'' and ``factored
\textbf{completely}.'' And carry the asymmetry with you: \textbf{one bad value disproves equivalence
forever, but no number of good values ever proves it} --- which is why the two most common errors of
the whole course, $(x+5)^2=x^2+25$ and $\sqrt{a+b}=\sqrt a+\sqrt b$, each die to a single number the
moment you test one.
\end{remindbox}

\end{document}
